%----------------------------------------------------------------------------------------
% General-purpose CV — DATA SCIENCE track (quick LinkedIn/Indeed applications)
% Reusable; not tied to a specific posting.
%----------------------------------------------------------------------------------------

\documentclass[10pt,a4paper,sans]{moderncv}

\usepackage[utf8]{inputenc}
\usepackage[T1]{fontenc}
\usepackage{lmodern}

\moderncvstyle{classic}
\moderncvcolor{blue}

\usepackage[scale=0.78]{geometry}
\usepackage{ragged2e}

%----------------------------------------------------------------------------------------

\firstname{Kundai Farai}
\familyname{Sachikonye}

\address{Zeltnerstrasse 30}{90443, Nuernberg,Germany}
\mobile{(+49) 1626386344}
\email{kundai.sachikonye@bitspark.com}
\homepage{github.com/fullscreen-triangle}
\photo[70pt][0.4pt]{pictures/Bew2.jpg}

\quote{Data scientist: machine learning, statistics, and reproducible pipelines
over large scientific datasets --- from raw data to deployed, validated tool.}

%----------------------------------------------------------------------------------------

\begin{document}

\makecvtitle

%----------------------------------------------------------------------------------------

\section{Profile}
\cvitem{}{
  Computational data scientist with deep experience in high-throughput biological
  data --- multi-omics, mass spectrometry, and imaging. I build statistically
  rigorous, reproducible machine-learning pipelines at scale, and take them to
  deployed tools validated against reference data.
}

\section{Technical Skills}
\cvitem{Analysis \& statistics}{
  Exploratory and statistical analysis; high-dimensional data; hypothesis testing;
  uncertainty quantification; experimental design
}
\cvitem{Machine learning}{
  PyTorch, scikit-learn, XGBoost; CNNs; representation learning; model evaluation
  and benchmarking
}
\cvitem{Data engineering}{
  Ray, Dask, SLURM/HPC; Docker, Kubernetes; HDF5, Parquet, memory-mapped I/O for
  $>$100\,GB datasets; reproducible pipelines
}
\cvitem{Programming}{
  Python (primary), R, SQL, Bash, Rust
}
\cvitem{Data domains}{
  Genomics / transcriptomics, proteomics / metabolomics (mass spectrometry),
  biomedical imaging
}

\section{Experience}
\cventry{2021--present}{Independent Data Science \& Computational Biology Researcher}{Bitspark GmbH}{}{}{
  Full-time development of ML pipelines over multi-omics, mass-spectrometry, and
  imaging data; deployed, documented, reference-validated tools.
}
\cventry{2019--2021}{Computational Lipidomics --- Data Analysis \& HPC}{Technische Universität München}{Munich}{}{
  High-throughput analysis at HPC scale; statistical evaluation of $>$100\,GB
  datasets; PyTorch classification. Taught at Master's level.
}
\cventry{2017--2018}{Glycomics \& Proteomics}{Universitätsklinikum Hamburg-Eppendorf}{Hamburg}{}{
  Automated LC-MS/MS and MALDI-TOF data pipelines; reproducible workflows.
}
\cventry{2013}{Cancer Biomarker Discovery}{University Medical Center Groningen}{Groningen}{}{
  UPLC-MS/MS proteomics; analysis of high-dimensional feature sets.
}

\section{Selected Projects}
\cvitem{gospel}{
  Multi-omics integration (genotype + transcriptomics + metagenomics); validated
  against 1000~Genomes, gnomAD, UK~Biobank.
  \href{https://github.com/fullscreen-triangle/gospel}{github.com/fullscreen-triangle/gospel}
}
\cvitem{lavoisier}{
  High-throughput mass-spectrometry pipeline; 98.9\% NIST accuracy; distributed
  (Ray). \href{https://lavoisier.vercel.app/sandbox}{lavoisier.vercel.app/sandbox}
}
\cvitem{syndrome}{
  Computational oncology --- neoantigen / MHC prediction (AUC~0.81, IEDB).
  \href{https://syndrome-seven.vercel.app/tools}{syndrome-seven.vercel.app}
}

\section{Education}
\cventry{2018--2019}{MSc Computational Biology}{Universität Tübingen}{}{}{}
\cventry{2014--2017}{MSc Pharmaceutical Biotechnology}{HAW Hamburg}{}{}{}
\cventry{2011--2014}{BSc Biochemistry and Cell Biology}{Jacobs University Bremen}{}{}{}
\cvitem{}{\textit{Additional:} doctoral research, computational lipidomics, TUM (not completed).}

\section{Languages}
\cvitemwithcomment{English}{Native / Fluent (C2)}{}
\cvitemwithcomment{German}{Conversational}{resident in Germany since 2011}

\renewcommand{\listitemsymbol}{-~}

\end{document}
